\section{Síntesis y ampliación}

Esta clase introductoria cumplió tres objetivos: definir la estructura del sistema de un Agente de LLM, explicar por qué el framework de Agentes merece estudiarse por separado, y presentar el mapa de capacidades y la hoja de ruta de proyectos de todo el curso. Su mensaje central es muy directo: las capacidades de los modelos grandes ya son suficientemente sólidas, pero para abordar tareas del mundo real el cuello de botella clave pasa de «si el modelo sabe decirlo» a «si el sistema sabe hacerlo».

\subsection{Lecturas adicionales}

\begin{itemize}
    \item Mialon et al., ``GAIA: a benchmark for General AI Assistants,'' 2023--2024.
    \item Zhou et al., ``WebArena: A Realistic Web Environment for Building Autonomous Agents,'' 2024.
    \item Jimenez et al., ``SWE-Bench: Can Language Models Resolve Real-World GitHub Issues?,'' 2024.
    \item Yao et al., ``ReAct: Synergizing Reasoning and Acting in Language Models,'' ICLR 2023.
\end{itemize}


\end{document}
